% Pursuing Stacks (cotes 134-2 to 134-8): archive cote 134-5, site page 118 (PDF page 119), typescript page 305.
% typescript. Draft compared: ps5.tex:668-738.
% AI-transcribed by Claude Opus 5.5 (claude-opus-5-5, adaptive thinking, effort high), Message Batches API, batch ps3 (msgbatch_01XdFmuVVBAcgc7EFxdtCqsj), 2026-11-12; reviewed against both supplied scan halves; see CHECK_DECISIONS.tsv.
\scanpage{134-5/118}
\mspage{305}
We should of course define composition of the \hand{\uline{\textcolor{black}{Hom}}}(A)'s ,  as required
too in Quillen's  set up.  This is done by relating the \hand{\uline{\textcolor{black}{Hom}}}(A)'s  to
the well-known interhal  \hand{\uline{\textcolor{black}{Hom}}}'s  in $M^{\wedge}$  -  which will show at the same
time that \deleted{in}\inserted[typed]{for} formula (3) to make sense,  we do not really have to assume
that M be stable under binary products,  as we can interprete the pro-%
ducts  $i(a)\mathrm{x}\hand{x}$  as being taken in $M^{\wedge}$,  as we\over{\unsure{ll}}{ll} as the Hom ,  so as to
get
\[ \mathrm{Hom}_{M^{\wedge}}(i(a)\mathrm{x}\hand{x} \, , \, y) \hand{\simeq} \mathrm{Hom}_{M^{\wedge}}(i(a), \, \hand{\underline{\textcolor{black}{\mathrm{Hom}}}}_{M^{\wedge}}(x,y)) \quad , \]
hence
\[ (4) \qquad \hand{\underline{\textcolor{black}{\mathrm{Hom}}}}(A)(x,y) \hand{\simeq} i^{*}(\hand{\underline{\textcolor{black}{\mathrm{Hom}}}}_{M^{\wedge}}(x,y)) \quad , \]
where $i^{*}$  in the right hand side is interpreted as a functor
\[ i^{*} \ : \ M^{\wedge} \hand{\longrightarrow} A^{\wedge} \hand{,} \]
rather than $M \hand{\longrightarrow} A^{\wedge}$ .  (I leave to the reader the task of enlarging the
basic Universe,  as need may be ...)  As $i^{*}$ commutes to products,  the
evident composition of t\over{\unsure{h}}{h}e internal \hand{\uline{\textcolor{black}{Hom}}}'s  in $M^{\wedge}$ gives rise to the
looked-for composition of the \hand{\uline{\textcolor{black}{Hom}}}(A)'s,  with the required associativity
properties.  Of course,  in case M is stable under binary products and
\hand{\uline{\textcolor{black}{Hom}}}'s,  \xout{in the} which apparently is going to be the case in all modeli-%
zing situations,  there is no need in the interpretation (4) to intro-%
duce the prohibitively large $M^{\wedge}$,  and we can work in M throughout.

There is an important relation though on the external \hand{\uline{\textcolor{black}{Hom}}}(A)'s
which we would like to be true for a satisfactory formalism,  namely
\[ (5) \qquad \hand{\Gamma}_{A^{\wedge}}(\hand{\underline{\textcolor{black}{\mathrm{Hom}}}}(A)(x,y)) \hand{\xleftarrow{\ \sim\ }} \mathrm{Hom}_{M}(x,y) \quad , \]
where $\hand{\Gamma}_{A^{\wedge}}$ just means  $\mathrm{Hom}_{A^{\wedge}}(e_{A^{\wedge}} \, , \, \ldots)$ .   This is equivalent  to
the requirement
\[ (6) \qquad i_{!}(e_{A^{\wedge}} \, ) \hand{\simeq} e_{M} \]
(assuming a final object $e_{M}$ \xout{\unsure{)}}in M to exist),  or equivalently
\[ (7) \qquad i(e_{A}) = e_{M} \]
if we assume moreover $e_{A}$ to exist.   Thus,  it will be natural to \xout{try}
ask for test functors satisfying the extra condition (5) or (6)  -
and when trying to construct test functors in various situations
(even without being aware of constructing test functors,  as Mr Jour-%
dain was ``doing prose without knowing it'' ...),  the very first thing
everybody has been doing instinctively was to write down formula (7),

\margin{118}
